\section{Deep discrepancy and interaction estimates}
\label{sec:the-two-deep-orientations-some-estimates}

\begin{corollary}[Deep discrepancy in both orientations]\label{cor:deep-discrepancy}
In the remaining case, on the simultaneous retained sets with the full cluster exclusions, there are constants \(0<x_*<.01,\ 0<\alpha<.01\) for which
\begin{equation}
 |d_G(\sigma,\pi)-1/2|\le b_*:=n^{-1+.04}
 \quad\text{at widths } \alpha n,\ n^{x_*},\text{ in either order}. \label{eq:source-9}
\end{equation}
This is for either color, including arbitrary probability laws at the indicated widths. The analogous assertion for any fixed smaller positive exponent slack instead of \(.04\) uses possibly smaller positive budget constants.
\end{corollary}

\begin{proof}
Apply Corollary~\ref{cor:remaining-deep-regime} with the chosen positive error slack, and shrink the two budget constants to work simultaneously in both orientations.
\end{proof}

We next bound the lower tail of common-neighbor mass after independent second-side draws. Small interactions are controlled by centered moments; in the homogeneous case, clique exclusion also bounds the number of unusually correlated labels.

Below, either orientation of the retained sides can be denoted \(X,Y\). Let \(f_x(y)=2\mathbf 1_G(x,y)-1\). Write \(m_\pi(x)=\int f_x\,d\pi,\ D_\pi(x)=(1+m_\pi(x))/2,\ K_\pi(x,z)=\int f_x f_z\,d\pi\).

Here and later, constants for small failure probabilities can be fixed in the following order. Take a cell-size exponent \(A_c=200\), a sufficiently large constant integer \(P\) (arbitrarily large in terms of \(A_c\), and of the absolute constants and polynomial exponents for the clock lemma with \(B=5\)), then \(R=P^2\). We describe estimates with this much slack even where a smaller power suffices. Call laws of widths \(\le n^{x_*/4}\) small-width in this section. We use this terminology only where later fixed multiples or moderate conditionings still fit \eqref{eq:source-9}.

\begin{lemma}[Heterogeneous centered-moment identity]\label{lem:heterogeneous-centered-moment}
Use \(d\le n\) independent inputs with possibly different small-width laws \(\pi_l\) on \(Y\), and a small-width probability \(\tau\) on \(X\) with \(|D_{\pi_l}-.5|\le O(b_*)\) on its support for each \(l\). Write \(a_{l,x}= (f_x-m_{\pi_l}(x))/(1+m_{\pi_l}(x))\), so \(1+a_{l,x}=\mathbf1_G(x,\cdot)/D_{\pi_l}(x)\). Put
\[
 Z(\boldsymbol y)=\int\prod_{l\le d}(1+a_{l,x}(y_l))\,d\tau(x).
\]
For even \(u\), the expectation \(\mathbb E_{\boldsymbol y}(Z-1)^u\), over the independent \(y_l\sim\pi_l\), is the integral for \(x_1,\ldots,x_u\) iid from \(\tau\) of
\begin{equation}
 \sum_{I\subseteq[u]}(-1)^{u-|I|}
       \prod_{l\le d}\int \prod_{i\in I}(1+a_{l,x_i})\,d\pi_l .
 \label{eq:source-10}
\end{equation}
Interactions are \(M_{l,J}=\int\prod_{i\in J}a_{l,x_i}\,d\pi_l,\ |J|\ge2\) (an empty maximum of interaction magnitudes is 0).
\end{lemma}

\begin{proof}
Expand the centered power, introduce independent copies of the first label for its factors, and integrate the independent second labels one coordinate at a time. The alternating sum is exactly \eqref{eq:source-10}; the individual centered factors have zero mean, so singleton interactions vanish.
\end{proof}

\begin{lemma}[Interaction tails]\label{lem:deep-interaction-tails}
In the setting of Lemma~\ref{lem:heterogeneous-centered-moment}, the following estimates hold for each fixed \(u\). A free label is an unfixed tuple coordinate, sampled with its original law \(\tau\); the other tuple coordinates are held fixed.
\begin{itemize}
\item With one free label and the other labels fixed, the magnitude of any interaction involving it exceeds \(O_u(b_*)\) with probability at most \(e^{-\alpha n/3}\).
\item With two involved labels free and the others fixed, it exceeds \(n^{-1-.03}\) with probability at most \(e^{-n^{.4}}\). The same two assertions hold for \(K_{\pi_l}\) on the labels at a pair of distinct tuple coordinates, even without the degree assumption in that raw test.
\item \(\mathbb E |M_{l,J}|\le n^{-.4|J|}\).
\item The one-free bound also holds with the degree requirements retained as a gate: if the other coordinates satisfy them, the unnormalized probability that the free coordinate satisfies them and has interaction magnitude above \(O_u(b_*)\) is at most \(e^{-\alpha n/3}\). Here the incoming support need not satisfy the degree requirements everywhere; the interaction is used only on the gate.
\end{itemize}
\end{lemma}

\begin{proof}\leavevmode\par
\paragraph{One free coordinate.}
Fix the column \(l\) and the other coordinates. The interaction is the projection of \(a_{l,x}\) onto a fixed bounded function of the second label. Equation~\eqref{eq:source-9} applies to bounded signed tests: add a fixed positive offset, normalize to a probability law, apply discrepancy, and subtract the offset term. If a one-sign exceptional set of first labels had \(\tau\)-mass at least \(\tfrac12e^{-\alpha n/3}\), conditioning on it would have width
\[
 n^{x_*/4}+\alpha n/3+O(1)<\alpha n.
\]
The signed-test estimate would then bound the alleged exceptional projection by \(O_u(b_*)\), a contradiction. This also proves the gated assertion: condition on the signed exception intersected with the degree gate, and never divide by the probability of the gate alone.

\paragraph{Two free coordinates.}
Write the interaction as \(\langle a_{l,x},H a_{l,z}\rangle_{\pi_l}\), with \(H\) a fixed bounded function supplied by the other coordinates. Suppose its magnitude exceeds \(n^{-1-.03}\) on a set of \(\tau^2\)-mass greater than \(e^{-n^{.4}}\). Splitting signs and then the conditional masses gives a set \(A\) of first coordinates, of mass at least \(e^{-O(n^{.4})}\), such that for every \(x\in A\) the same signed violation holds on a set \(E_x\) of \(z\)'s with \(\tau(E_x)\ge e^{-O(n^{.4})}\).

Let \(\rho_x=\tau(\cdot\mid E_x)\). Its width is \(O(n^{.4})<\alpha n\). Applying \eqref{eq:source-9} to this first law shows that \(\mathbb E_{\rho_x}f_z(y)=O(b_*)\) outside a \(\pi_l\)-set of mass \(e^{-\Omega(n^{x_*})}\). The degree bounds change \(f_z\) to \(a_{l,z}\) by \(O(b_*)\), and boundedness on the exceptional set therefore gives
\(\|\mathbb E_{\rho_x}a_{l,z}\|_{L^2(\pi_l)}=O(b_*)\).
Take \(t=\lceil n^{.25}\rceil\) independent \(z_j\sim\rho_x\). For \(S=H\sum_{j\le t}a_{l,z_j}\), independence gives
\[
 \mathbb E\|S\|_2^2=O_u(t+t^2b_*^2)=O_u(t).
\]
A fixed positive fraction of these tuples consequently satisfy \(\|S\|_2=O_u(\sqrt t)\).

Return now to the original product law \(\tau^{\otimes t}\). For each \(x\in A\), the tuples just obtained have raw mass at least \(e^{-O(tn^{.4})}\), and all their coordinates lie in \(E_x\). Fubini therefore supplies one fixed norm-good tuple for which the set of \(x\)'s satisfying all the signed violations has mass at least \(e^{-O(tn^{.4})}\). Conditioning \(\tau\) on that set costs width \(O(n^{.65})<\alpha n\). Its mean \(a_{l,x}\) has \(L^2(\pi_l)\)-norm \(O(b_*)\) by the preceding argument. Its projection onto this fixed \(S\) is thus at most
\[
 O_u(b_*\sqrt t)=O_u(n^{-.835}),
\]
whereas the simultaneous violations make that projection at least
\(t n^{-1-.03}=n^{-.78+o(1)}\) with one sign. This contradiction proves the two-free estimate. For the raw correlation \(K_{\pi_l}(x,z)=\langle f_x,f_z\rangle_{\pi_l}\), the same proof uses \(f\) in place of \(a\) and requires no degree gate.

\paragraph{Mean interaction size.}
Independence of the tuple coordinates yields
\[
 \mathbb E M_{l,J}^2
 =\mathbb E_{y,y'\sim\pi_l}
   \left(\int a_{l,x}(y)a_{l,x}(y')\,d\tau(x)\right)^{|J|}.
\]
For fixed \(y\), the inner expression is a bounded signed test on the first side applied to \(f_x(y')\), with an \(O(b_*)\) error from the centering and denominators. Deep discrepancy bounds it by \(O(b_*)\) outside a \(\pi_l\)-set of \(y'\)'s of mass \(e^{-\Omega(n^{x_*})}\). Boundedness elsewhere gives
\[
 \mathbb E M_{l,J}^2\le (C_u b_*)^{|J|}
             +C_u e^{-\Omega(n^{x_*})}.
\]
Cauchy--Schwarz now gives \(\mathbb E|M_{l,J}|\le n^{-.4|J|}\), since \(.96/2>.4\).
\end{proof}

Put \(L=\lceil500R\rceil\), and choose an even constant \(u>10L^2\). Take successively tiny constants \(\xi,\theta>0\), with \(\xi<\alpha 2^{-10u-100}\) and \(\theta\) sufficiently small in terms of \(\xi,u\). All can be fixed before patch and geometry choices below.

\begin{lemma}[The moderate-interaction contribution]\label{lem:moderate-interaction-moment}
Let \(\Phi_u\) be the signed integrand in \eqref{eq:source-10}, and put \(w=\max_{l,J}|M_{l,J}|\). Then
\[
 \left|\int_{\{w\le2\xi\}}\Phi_u\,d\tau^{\otimes u}\right|\le n^{-3R}.
\]
Each positive product term in that integrand has bounded integral on the same set. This latter assertion also holds for subtuples of the fixed length bound.
\end{lemma}

\begin{proof}
Every positive product term is at most \(\exp(2^u n w)\). The union over the \(O_u(n)\) interactions and Lemma~\ref{lem:deep-interaction-tails} bound its integral on the two ranges above \(n^{-1-.03}\) by
\[
 \begin{array}{ll}
 O_u(n)e^{2^u n^{.06}-n^{.4}},
   & n^{-1-.03}<w\le n^{-1+.06},\\
 O_u(n)e^{2^{u+1}\xi n-\alpha n/3},
   & n^{-1+.06}<w\le2\xi.
 \end{array}
\]
In the second range the lower endpoint exceeds the one-free threshold \(O_u(b_*)\). Both bounds decay faster than any fixed inverse power of \(n\); the second uses the stipulated smallness of \(\xi\).

On \(w\le n^{-1-.03}\), expand each column factor into interactions and then sum over \(I\). A term with chosen interactions whose union is \(S\subseteq[u]\) has coefficient \(\sum_{I\supseteq S}(-1)^{u-|I|}\). It vanishes unless \(S=[u]\). Thus every surviving term covers all \(u\) coordinates, and its interactions come from distinct columns. Terms choosing more than \(L\) interactions have a geometric tail bounded by powers of \(2^u n^{-.03}\), which is smaller than the required \(n^{-3R}\). Among at most \(L\) interactions covering \([u]\), one has at least \(u/L\) coordinates. Its mean absolute size is at most \(n^{-.4u/L}\); the remaining factors are bounded on this set. The column and interaction choices cost \(O_u(n^L)\), giving \(O_u(n^{L-.4u/L})\), again smaller than \(n^{-3R}\) because \(u>10L^2\).

Without the alternating sum, the small range has weight \(\exp(O_u(n^{-.03}))\), and the two displayed tail bounds still apply. This proves the uncentered assertion, including for subtuples.
\end{proof}

\begin{lemma}[Homogeneous extension counts and peeling]\label{lem:homogeneous-peeling}
Use the homogeneous version \(\pi_l=\pi\). Suppose the support in use has no clique on \(\lceil e^Q\rceil\) labels with \(K_\pi>3\theta\) between all distinct members. Then given any other labels of a tuple, the number of extensions by one label from the support to any interaction of magnitude \(>\xi\) involving that label is at most \(\exp(C_* Q)\), for a fixed \(C_*\) (we use \(Q\) large). Also for any fixed \(x\), the number of labels \(z\) there with \(|K_\pi(x,z)|>\xi\) satisfies this bound. If in addition
\[
 \max_x\tau(x)D_\pi(x)^{-d}\exp(C_*Q)\le\Gamma<1
\]
on the support, then the absolute contribution of the large-interaction tuples to the positive product terms of \eqref{eq:source-10}, summed over those terms, is \(O_u(\Gamma)\).
\end{lemma}

\begin{proof}
After the other coordinates are fixed, a large extension gives a projection of \(f_z\) onto a bounded fixed direction whose magnitude is bounded below by a positive constant depending only on \(\xi,u\). Split the candidates by projection sign. If \(t\) vectors in one sign class have all mutual inner products at most \(3\theta\), their average has squared norm at most \(1/t+3\theta\), while its projection has that fixed positive lower bound. For sufficiently small \(\theta\), this forces \(t<t_0\) for a fixed \(t_0=t_0(\xi,u)\). The graph joining pairs with \(K_\pi>3\theta\) has neither an independent set of size \(t_0\) nor a clique of size \(\lceil e^Q\rceil\). The elementary Ramsey binomial bound gives
\[
 \log\binom{\lceil e^Q\rceil+t_0}{t_0}\le C_*Q.
\]
Increase the fixed \(C_*\), chosen after \(u,\xi,\theta\), to include the two signs and the finitely many interactions. The same projection argument gives the stated count for \(|K_\pi(x,z)|>\xi\).

For a tuple with a large interaction, choose a maximal retained index set whose subtuple has every interaction at most \(\xi\). Sum over the finitely many possible retained index sets. Once its labels are fixed, each excluded coordinate must create a large interaction when added to them; otherwise maximality fails. Its possible labels therefore lie in a set of size at most \(\exp(C_*Q)\), determined by the retained labels alone.

In a positive product term of \eqref{eq:source-10}, deleting an excluded coordinate costs at most \(D_\pi(x)^{-d}\). Under the original product measure, each excluded coordinate consequently contributes at most
\[
 \exp(C_*Q)\max_x\tau(x)D_\pi(x)^{-d}\le\Gamma.
\]
These necessary extension conditions can be integrated independently; no independence after conditioning on maximality is asserted. The remaining positive product has bounded integral by Lemma~\ref{lem:moderate-interaction-moment}. Every tuple under consideration excludes at least one coordinate, so summing the finitely many retained sets and positive terms gives \(C_u\Gamma\).

In particular, for the homogeneous mass \(Z\) defined in Lemma~\ref{lem:heterogeneous-centered-moment}, the moderate bound and even centered-power Markov give the explicit consequence
\[
 \Pr\{Z<1/2\}\le2^u\bigl(n^{-3R}+C_u\Gamma\bigr).
\]
The peeling bound on the set with an interaction above \(\xi\) covers the complement of the moderate set. In heterogeneous applications the large-interaction part will be controlled separately.
\end{proof}

The small- and moderate-interaction estimates of Lemmas~\ref{lem:deep-interaction-tails} and~\ref{lem:moderate-interaction-moment} allow a different law at each second-label coordinate. The clique-count argument of Lemma~\ref{lem:homogeneous-peeling} is used only with a single common law. In particular, conditioning on internal data later does not convert a heterogeneous collection into a homogeneous one.

\begin{lemma}[Simultaneous row trimming]\label{lem:simultaneous-row-trimming}
For a uniform small-width law on a set of first labels, and any collection of at most \(\exp(o(n^{.2}))\) small-width second laws, we can, by removing an \(o(1)\) fraction, arrange (for each law) for every kept \(x\), with \(z\) uniform on the original first set:
\begin{equation}
 \Pr_z(|K_\pi(x,z)|>n^{-1-.03})\le e^{-n^{.2}},\qquad
 \mathbb E_z\big[\mathbf1_{\{n^{-1-.03}<|K_\pi(x,z)|\le4\xi\}}
                       e^{150n|K_\pi(x,z)|}\big]\le e^{-n^{.2}} .
 \label{eq:source-11}
\end{equation}
For parameters subsequently interpolated at extremely fine mesh we can use \(n^{-1-.02},\,2\xi,\,100n,\,e^{-n^{.19}}\) instead at the realized laws, by choosing the mesh fine enough.
\end{lemma}

\begin{proof}
Let \(x,z\) be independent uniform labels in the original first set. The raw two-free estimate gives
\(\Pr(|K_\pi(x,z)|>n^{-1-.03})\le e^{-n^{.4}}\).
For the weighted expression, split at \(n^{-1+.06}\). The two-free and one-free bounds respectively give
\[
 \mathbb E_{x,z}\!\left[
  \mathbf1_{\{n^{-1-.03}<|K_\pi(x,z)|\le4\xi\}}
        e^{150n|K_\pi(x,z)|}\right]
 \le e^{150n^{.06}-n^{.4}}+e^{600\xi n-\alpha n/3}
 \le e^{-c n^{.4}}
\]
for a fixed \(c>0\) and large \(n\). Markov, applied to each conditional expectation over \(z\) at threshold \(e^{-n^{.2}}\), removes at most \(e^{-c n^{.4}+n^{.2}}\) of the first labels per test. The union over \(\exp(o(n^{.2}))\) laws still removes an \(o(1)\) fraction.

For interpolation, choose simplex diameters at most \(n^{-3}\) in \(\ell^1\). Then \(|K_\pi-K_{\pi'}|\le\|\pi-\pi'\|_1\le n^{-3}\). The realized event with thresholds \(n^{-1-.02},2\xi\) is contained in the corner event with thresholds \(n^{-1-.03},4\xi\). Its exponential weight is at most the old weight times \(e^{100n^{-2}}\). The relaxation from \(e^{-n^{.2}}\) to \(e^{-n^{.19}}\) pays for this factor. A finer mesh can simultaneously preserve any other finitely many strict margins needed later.
\end{proof}

\paragraph{Parameters and normalizations in the final regime.}
The constants and scales in Sections~\ref{sec:the-two-deep-orientations-some-estimates}--\ref{sec:late-dynamics-and-its-transfer-problem} have the following separate roles.
\begin{itemize}
\item The failure-probability constants are fixed as
\(A_c=200\), then sufficiently large \(P\), then \(R=P^2\).
The moment estimates use a sufficiently large even \(u\) depending on
\(R\), followed by small \(\xi\) and then \(\theta\), as specified above.
Their homogeneous extension bound supplies the fixed constant \(C_*\).
\item Section~\ref{sec:measured-patch-scales} chooses width exponents
within the available discrepancy budgets.
Section~\ref{sec:extraction-and-cube-allocation} fixes the extraction and
geometry constants in the order required by the construction, ending with
a sufficiently large fixed scale threshold. Estimates above that threshold
are uniform; low cluster scales need not tend to infinity.
\item Extraction fixes residual bias and cluster scales, patch sizes,
internal dimensions, and gains. Mode, orientation, and color selection
precedes prefix allocation and the later choice of internal coordinate
sets. These quantities may depend on \(n\). The low regime uses order \(\log n\)
late classes and \(\lceil\log^2 n\rceil\) resampling rounds. Its late-error
exponent is fixed small after the low-mode constants, as specified in
Section~\ref{sec:late-dynamics-and-its-transfer-problem}.
\item Widths remain normalized by the original \(N\), including after
labels are discarded. The later local normalizations are introduced where
used: patch loads refer to the extracted side size, endpoint kernels to that
size divided by the number of palettes, and late kernels to the size of the
current late pool.
Temporary notation has the scope declared at its use.
\end{itemize}
